%% starter: A beamer talk on a lab: title, motivation, method, results and conclusion slides
\documentclass{beamer}
\usetheme{default}
\usepackage{amsmath}
\usepackage{graphicx}
\usepackage{siunitx}

\title{{{Title}}}
\author{{{Author}}}
\institute{Department of Physics}
\date{\today}

\begin{document}

% Each frame is one slide; \frametitle names it.
\begin{frame}
  \titlepage
\end{frame}

\begin{frame}
  \frametitle{Motivation}
  % Keep a slide to a few short points; say the rest aloud.
  \begin{itemize}
    \item Why the experiment is worth doing
    \item What it sets out to measure
  \end{itemize}
\end{frame}

\begin{frame}
  \frametitle{Method}
  The period of a pendulum of length $L$ is
  \begin{equation*}
    T = 2\pi\sqrt{\frac{L}{g}},
  \end{equation*}
  so we timed ten swings at each length from \qty{0.20}{m} to \qty{1.00}{m}.
\end{frame}

\begin{frame}
  \frametitle{Results}
  \begin{figure}
    % Put your plot in figures/ and uncomment the next line, naming it.
    % \includegraphics[width=0.8\linewidth]{figures/plot.pdf}
    \fbox{\parbox[c][4cm][c]{0.8\linewidth}{\centering Your figure here}}
    \caption{The period $T$ against $\sqrt{L}$, with a straight-line fit.}
  \end{figure}
\end{frame}

\begin{frame}
  \frametitle{Conclusion}
  \begin{itemize}
    \item What the results show
    \item How they compare with the accepted value
  \end{itemize}
\end{frame}

\end{document}
